\section{Traditional, ANN, and Model2 Methods}
\label{sec:methods-comparison}

\subsection{Model 1: ordinary ANN}
Model 1 consumes the frozen 100-feature R2 vector and predicts the ten canonical
parameters through a hard output map. It uses seeds 11, 22, and 33. Validation-only
checkpoint selection is recorded for each seed; the test split is not used for
selection.

\subsection{Model 2: constraint- and repricing-informed inverse model}
Model 2 adds structural output enforcement and a differentiable Fourier-repricing
loss to the inverse mapping. It is not PDE-informed. Its primary cohort comprises the
same three seeds trained uniformly on Tesla P100 hardware. A local CPU replication of
seed 11 is retained as execution provenance but excluded from primary inference.

\subsection{Traditional numerical calibration}
The traditional baseline uses the frozen production pricer, trust-region reflective
least squares, three starts per surface, start seed 42, at most 300 function
evaluations per start, tight tolerances, and reviewed provisional bounds. The frozen
representative rule selects the start with lowest final objective, breaking ties by
lowest start index. Every start, including failures and budget stops, remains in the
journal.

The traditional run received a disclosed truth-informed start favoring it by protocol.
This makes it a strong accuracy reference rather than a claim that arbitrary starts
identify the truth. Runtime and training environments are not compared as scientific
quality.
